% Generated by experiments/scripts/export_latex.py -- do not edit by hand.
% Source: experiments/data/2-processed-tables-and-evaluations/revised-usd/same-window/eval_summary.json
% Source: experiments/data/2-processed-tables-and-evaluations/revised-usd/holdout-w0/eval_summary.json
% Source: experiments/data/2-processed-tables-and-evaluations/revised-usd/window-w1/eval_summary.json
\begin{table}[htbp]
\centering
\caption[AWP 옆에 놓은, 고르게 다시 출발하는 엔도스랭크와 활동에 따라 무게를 단 재시작을 쓰는 엔도스랭크.]{엔도스랭크를 고르게 다시 출발하게 한 것과, 허락하는 활동에 따라 무게를 단 재시작으로 다시 출발하게 한 것을 AWP 옆에 놓았어요. 뒤의 재시작은 AWP의 규칙을 허락 그래프에 쓴 것이에요. 매칭 코호트(계약 $n=15{,}107$개)의 지표 가족 평균 켄달의 $\tau$를 보여 줘요. 또 두 스펜더 라벨과의 $\tau$를 W0에서(스펜더 계약 $n=17{,}050$개, 점수는 2026년 3월 31일에 동결), 그리고 W1에서($n=18{,}152$, 2026년 6월 30일에 동결) 보여 줘요. 두 엔도스랭크 행 모두 각 마지막 허락에 $\sigma(\Delta t)V(z)$만큼 무게를 줘요. 신뢰구간(CI)은 계약을 다시 뽑아 구한 95\% 백분위 부트스트랩 구간이에요(짝지은 다시 뽑기 400번, 시드 42).}
\label{tab:endorserank-restarts}
\fitwidth{%
\begin{tabular}{llccccccc}
\toprule
 & & \multicolumn{3}{c}{같은 기간, 매칭 코호트} & \multicolumn{2}{c}{W0, 스펜더} & \multicolumn{2}{c}{W1, 스펜더} \\
\cmidrule(lr){3-5} \cmidrule(lr){6-7} \cmidrule(lr){8-9}
방법 & 재시작 & 송금 & 허락 & 시빌 & 새 허락 쌍 & 새 송금자 & 새 허락 쌍 & 새 송금자 \\
\midrule
EndorseRank & 고르게 & \begin{tabular}[c]{@{}c@{}}0.158\\[-1pt]{\footnotesize [0.147, 0.169]}\end{tabular} & \begin{tabular}[c]{@{}c@{}}0.585\\[-1pt]{\footnotesize [0.578, 0.592]}\end{tabular} & \begin{tabular}[c]{@{}c@{}}0.007\\[-1pt]{\footnotesize [-0.001, 0.015]}\end{tabular} & \begin{tabular}[c]{@{}c@{}}0.231\\[-1pt]{\footnotesize [0.220, 0.242]}\end{tabular} & \begin{tabular}[c]{@{}c@{}}0.200\\[-1pt]{\footnotesize [0.189, 0.210]}\end{tabular} & \begin{tabular}[c]{@{}c@{}}0.234\\[-1pt]{\footnotesize [0.222, 0.244]}\end{tabular} & \begin{tabular}[c]{@{}c@{}}0.194\\[-1pt]{\footnotesize [0.184, 0.205]}\end{tabular} \\
EndorseRank & 허락하는 활동에 따라 & \begin{tabular}[c]{@{}c@{}}0.164\\[-1pt]{\footnotesize [0.153, 0.173]}\end{tabular} & \begin{tabular}[c]{@{}c@{}}0.426\\[-1pt]{\footnotesize [0.418, 0.435]}\end{tabular} & \begin{tabular}[c]{@{}c@{}}0.049\\[-1pt]{\footnotesize [0.039, 0.057]}\end{tabular} & \begin{tabular}[c]{@{}c@{}}0.308\\[-1pt]{\footnotesize [0.299, 0.319]}\end{tabular} & \begin{tabular}[c]{@{}c@{}}0.269\\[-1pt]{\footnotesize [0.260, 0.278]}\end{tabular} & \begin{tabular}[c]{@{}c@{}}0.295\\[-1pt]{\footnotesize [0.285, 0.304]}\end{tabular} & \begin{tabular}[c]{@{}c@{}}0.242\\[-1pt]{\footnotesize [0.232, 0.253]}\end{tabular} \\
AWP & 송금하는 활동에 따라 & \begin{tabular}[c]{@{}c@{}}0.572\\[-1pt]{\footnotesize [0.564, 0.580]}\end{tabular} & \begin{tabular}[c]{@{}c@{}}0.193\\[-1pt]{\footnotesize [0.182, 0.203]}\end{tabular} & \begin{tabular}[c]{@{}c@{}}0.395\\[-1pt]{\footnotesize [0.387, 0.403]}\end{tabular} & \begin{tabular}[c]{@{}c@{}}0.196\\[-1pt]{\footnotesize [0.181, 0.210]}\end{tabular} & \begin{tabular}[c]{@{}c@{}}0.346\\[-1pt]{\footnotesize [0.336, 0.356]}\end{tabular} & \begin{tabular}[c]{@{}c@{}}0.169\\[-1pt]{\footnotesize [0.157, 0.183]}\end{tabular} & \begin{tabular}[c]{@{}c@{}}0.330\\[-1pt]{\footnotesize [0.321, 0.340]}\end{tabular} \\
\bottomrule
\end{tabular}
}
\end{table}
